\documentclass{beamer}

\usepackage[francais]{babel}
\usepackage{rtxslides}
\usepackage{tikz}
\usetikzlibrary{shapes,fit}

\title{\rtx, \emph{THE} device driver generator}
\date[Comité de Pilotage 3]{Comité de Pilotage 3 \\ \vspace{10pt} \includegraphics[height=12mm]{logo_eip}}
\author[\rtx\ 2012 -- 18 novembre 2011]{\rtx\ 2012 -- 18 novembre 2011 \vspace{-20pt}}

\hypersetup{pdfauthor={Rathaxes 2012},pdfsubject={Rathaxes, THE device driver generator}}

\definecolor{lightred}{RGB}{147,36,33}
\tikzset{componentarrow/.style={<-, >=latex, color=rathaxesred, ultra thick}}

\newcommand{\cemph}[1]{{\itshape{\textcolor{rathaxesred}{#1}}}}

\newcommand{\bgcolor}{black}
\newcommand{\fgcolor}{white}

\tikzset{redbox/.style={draw,rectangle,rounded corners=5pt,ultra thick,color=rathaxesred,text=\fgcolor}}
\tikzset{warrow/.style={->, >=stealth, color=\fgcolor, ultra thick}}

\setbeamertemplate{navigation symbols}{}

\begin{document}

\begin{frame}
\titlepage
\end{frame}

\begin{frame}{Rappel du Projet}
\rmfamily{\itshape{«~\rtx\ est un \textcolor{rathaxesred}{langage dédié} qui
permet de \textcolor{rathaxesred}{décrire des pilotes} de périphériques. \rtx\
compile vers des modules noyaux écrits en C pour \textcolor{rathaxesred}{Linux},
\textcolor{rathaxesred}{Windows} et \textcolor{rathaxesred}{OpenBSD}.~»}}
\end{frame}

\begin{frame}[fragile]{Rappel du Projet}
\begin{center}
\begin{tikzpicture}[overlay]
\draw (0, 2.7) node{\Large{S'adresse à deux types d'utilisateurs :}} ;

\draw (-2.5, 1.2) node[text width=2cm, text centered]{\small{Fabricants de Matériel}} ;
\draw (-2.5, -1.3) node[below]{\includegraphics[width=2cm]{logos/nvidia_white}} ;
\draw (-2.5, -0.3) node[below]{\includegraphics[width=2cm]{logos/intel}} ;
\draw (-2.5, 0.2) node[below]{\includegraphics[width=2cm]{logos/broadcom_white}} ;
\draw (-2.5, -3.3) node{\ldots} ;

\draw (2.5, 1.2) node[text width=2cm, text centered]{\small{Éditeurs de SE}} ;
\draw (2.5, 0.2) node[below]{\includegraphics[width=2cm]{logos/tux}} ;
\draw (2.5, -2.8) node[above]{\includegraphics[width=2cm]{logos/windows_7_white}} ;
\draw (2.5, -3.3) node{\ldots} ;
\end{tikzpicture}
\end{center}
\end{frame}

\begin{frame}[fragile]{Avancement}

\large{
\begin{itemize}
\item Version complète et utilisable du compilateur ;
\item Soutenances finales de 4\textsuperscript{ème} année positives ;
\item Supporte les trois parties du langage:
\end{itemize}
}

\begin{center}
\begin{tikzpicture}[overlay]
\draw (-3, -0.5) node{\includegraphics[width=1.5cm]{icons/rtx}} ;
\draw (0, -0.5) node{\includegraphics[width=1.5cm]{icons/rti}} ;
\draw (3, -0.5) node{\includegraphics[width=1.5cm]{icons/blt}} ;
\end{tikzpicture}
\end{center}
\end{frame}

\begin{frame}{Documentation}
\large{
\begin{itemize}
\item Mise à jour de la documentation technique ;
\item Mise à jour de la documentation utilisateur prévue.
\end{itemize}
}
\end{frame}

\begin{frame}[fragile]{Promotion}
\begin{center}
\begin{tikzpicture}
\node (RMLL) at (0, 0) {\includegraphics[height=6cm]{rmll_2011}};

\draw (RMLL.east) node[right] {\begin{minipage}{45mm}\begin{itemize}
\item \Large{RMLL 2011} ;
\item \Large{Youtube} ;
\item \Large{Google Code} ;
\item \Large{\emph{Forum EIP}} ;
\item \Large{\emph{Fosdem 2012}}.
\end{itemize}\end{minipage}};
\end{tikzpicture}
\end{center}
\end{frame}

\begin{frame}{Planning}
\Large{
\begin{itemize}
\item Pas de changements sur le CDC ;
\item Retard sur le compilateur ;
\item Soutenance technique finale :
\begin{itemize}
\Large{
\item Driver réseau PCI en \rtx.
}
\end{itemize}
\end{itemize}
}
\end{frame}

\begin{frame}{Merci}
\begin{center}
\Huge{\emph{Questions ?}}
\end{center}

\vspace{2em}
\begin{itemize}
\item \Large{\texttt{http://www.rathaxes.org/}}
\end{itemize}
\end{frame}

\end{document}
